\section{Intrinsic highest-multiplicity strata of normal-crossing carriers}
\label{sec:hironaka-normal-crossings-multiplicity-strata}

These statements concern the geometry of an original general
normal-crossing carrier. They produce an intrinsic first centre without
choosing an ordering of its local branches. They do not assert that this
centre is contained in the cosupport of an additional marked ideal.

\begin{definition}[Intrinsic local branch multiplicity]
\label{hir:general-nc-intrinsic-branch-multiplicity}
\lean{Hironaka.HironakaNormalCrossingsMultiplicityStrata.branchMultiplicity,Hironaka.HironakaNormalCrossingsMultiplicityStrata.branchMultiplicity_eq_card,Hironaka.HironakaNormalCrossingsMultiplicityStrata.branchMultiplicity_le,Hironaka.HironakaNormalCrossingsMultiplicityStrata.branchMultiplicity_pos_iff}
\uses{hir:actual-normal-crossing-chart-branch-correspondence}
\leanok
For an actual normal-crossing carrier $H$ in a complex manifold of
dimension $n$, let $b_H(x)$ be the number of distinct coordinate
hyperplanes in an actual centred flattening of its germ at $x$.
\end{definition}
\begin{proof}
An actual flattening exists by the normal-crossing condition. Any two
flattenings have the uniquely produced branch bijection of the overlap
theorem, so the numbers agree. Thus the definition is independent of
the flattening. Its values lie between zero and $n$. Since the
flattening sends $x$ to zero, $x$ lies in $H$ exactly when its selected
coordinate set is nonempty, equivalently when $b_H(x)>0$.
\end{proof}

\begin{lemma}[The exact nearby branch count]
\label{hir:general-nc-exact-nearby-branch-count}
\lean{Hironaka.HironakaNormalCrossingsMultiplicityStrata.branchMultiplicity_eq_filter_card}
\uses{hir:general-nc-intrinsic-branch-multiplicity}
\leanok
Fix an actual flattening $F$ at $x$ with selected coordinate set $S$.
On a neighbourhood of $x$, the original local branch multiplicity at
$y$ equals the number of coordinates $i\in S$ for which $F_i(y)=0$.
\end{lemma}
\begin{proof}
The flattening identifies $H$ with the selected hyperplane union on its
whole open domain. Fix a point $y$ in that domain and write
$S_y=\{i\in S:F_i(y)=0\}$. Every coordinate indexed by $S\setminus S_y$
is nonzero at $y$. The finitely many open nonvanishing neighbourhoods
have an intersection on which none of these coordinates contributes
to the union. Translate the actual coordinate map by $F(y)$. It has
an actual translated inverse and is centred at $y$. For $i\in S_y$,
its coordinate $i$ equals the original $F_i$, since $F_i(y)=0$.
Thus these coordinates describe exactly the original carrier germ at
$y$. Express this actual local holomorphic coordinate equivalence in
the preferred chart at $y$ using the two actual chart maps and their
inverses. It is an actual centred flattening with selected set $S_y$.
The defining overlap independence gives $b_H(y)=|S_y|$. The argument
holds for every $y$ in the domain of the original flattening, so that
whole domain is the required neighbourhood of $x$.
\end{proof}

\begin{lemma}[Closed intrinsic multiplicity superlevels]
\label{hir:general-nc-closed-multiplicity-superlevels}
\lean{Hironaka.HironakaNormalCrossingsMultiplicityStrata.isClosed_branchMultiplicity_superlevel,Hironaka.HironakaNormalCrossingsMultiplicityStrata.branchMultiplicity_superlevel_subset}
\uses{hir:general-nc-exact-nearby-branch-count}
\leanok
For every nonnegative integer $r$, the intrinsic set
$C_r=\{x:b_H(x)\ge r\}$ is closed. For $r>0$ it is contained in $H$.
\end{lemma}
\begin{proof}
At a point with $b_H(x)<r$, take an actual flattening with
$|S|=b_H(x)$. On its domain the preceding exact count is the cardinality
of a subset of $S$, so it remains strictly less than $r$. Thus the
complement of $C_r$ is open and $C_r$ is closed. If $r>0$, every
point of $C_r$ has positive branch multiplicity and hence lies in $H$.
\end{proof}

\begin{theorem}[The genuine highest-multiplicity centre]
\label{hir:general-nc-highest-multiplicity-smooth-centre}
\lean{Hironaka.HironakaNormalCrossingsMultiplicityStrata.highest_superlevel_flattening,Hironaka.HironakaNormalCrossingsMultiplicityStrata.regularAtDim_highest_superlevel,Hironaka.HironakaNormalCrossingsMultiplicityStrata.isClosedSubmanifold_highest_superlevel}
\uses{hir:general-nc-closed-multiplicity-superlevels,hir:general-nc-exact-nearby-branch-count,fpgeom:complex-analytic-subvariety}
\leanok
Suppose $r>0$ bounds the original branch multiplicity everywhere.
Then $C_r=\{x:b_H(x)=r\}$ is an actual closed smooth submanifold.
At each of its points it has dimension $n-r$, and an actual original
normal-crossing chart expresses it as the intersection of all of the
$r$ local original branches. These intersections agree as actual
germs on every chart overlap, independently of any branch ordering.
\end{theorem}
\begin{proof}
The upper bound identifies $C_r$ with the closed superlevel above.
At a point of $C_r$, take a flattening with selected set $S$ of size
$r$. The exact nearby-count theorem says a nearby point has count
at least $r$ exactly when every coordinate in $S$ vanishes. Thus the
actual centre germ is precisely the coordinate subspace cut out by
$S$. The original flattening has an invertible differential. Its distinct
selected coordinate rows are linearly independent: any linear relation
can be evaluated on the inverse image of each corresponding coordinate
basis vector to make its coefficient zero. These $r$ holomorphic
functions cut out the whole centre germ. The holomorphic implicit
function theorem therefore gives an actual regular flattening to the
dimension-$(n-r)$ coordinate plane; $r\le n$ follows from the
selected-set cardinality. Together with closedness this proves
the entire centre is an actual closed smooth submanifold.
The overlap theorem matches every local original branch uniquely;
its finite-intersection conclusion therefore identifies the whole
centre germ on every overlap. No matching or smooth-centre
certificate is assumed. If the level is empty, closed smoothness is
vacuous and the pointwise dimension assertion has no instances.
\end{proof}

\begin{corollary}[A nonempty intrinsic highest stratum]
\label{hir:general-nc-actual-highest-stratum-exists}
\lean{Hironaka.HironakaNormalCrossingsMultiplicityStrata.exists_highest_multiplicity_stratum}
\uses{hir:general-nc-highest-multiplicity-smooth-centre}
\leanok
If $H$ is nonempty, there is a positive integer $r\le n$ attained
by $b_H$ and bounding it everywhere, so its intrinsic highest
stratum is a nonempty actual closed smooth centre. No compactness
is required.
\end{corollary}
\begin{proof}
The nonempty range of $b_H$ is contained in the finite set
$\{0,\ldots,n\}$ and therefore has a largest attained value $r$.
A point of $H$ has positive branch multiplicity, hence $r>0$.
Choose a point attaining $r$; it lies in the highest stratum.
Apply the preceding theorem with the actual global upper bound.
In dimension zero the normal-crossing carrier must be empty, so
the nonemptiness hypothesis does not arise.
\end{proof}

\begin{remark}
An actual blowup of this intrinsic smooth centre is available from
the proved geometric existence theorem. Using it to lower the number
of strict original branches requires its genuine coordinate-chart
calculation and compatibility with any accumulated exceptional
boundary. A marked-ideal resolution that must preserve the complement
of its original cosupport cannot use this centre unless its required
cosupport containment has also been proved.
The ordered monomial invariant retains its colexicographic tie-break
on branch indices. The intrinsic branch count and its centre do not
prove that this ordered canonical sequence is independent of an
arbitrary initial permutation of those indices.
\end{remark}
